\documentclass[12pt,]{article}
\usepackage{lmodern}
\usepackage{amssymb,amsmath}
\usepackage{ifxetex,ifluatex}
\usepackage{fixltx2e} % provides \textsubscript
\ifnum 0\ifxetex 1\fi\ifluatex 1\fi=0 % if pdftex
  \usepackage[T1]{fontenc}
  \usepackage[utf8]{inputenc}
\else % if luatex or xelatex
  \ifxetex
    \usepackage{mathspec}
  \else
    \usepackage{fontspec}
  \fi
  \defaultfontfeatures{Ligatures=TeX,Scale=MatchLowercase}
\fi
% use upquote if available, for straight quotes in verbatim environments
\IfFileExists{upquote.sty}{\usepackage{upquote}}{}
% use microtype if available
\IfFileExists{microtype.sty}{%
\usepackage{microtype}
\UseMicrotypeSet[protrusion]{basicmath} % disable protrusion for tt fonts
}{}
\usepackage[margin=1in]{geometry}
\usepackage{hyperref}
\hypersetup{unicode=true,
            pdftitle={Exceedingly Simple Principal Pivot Transforms},
            pdfauthor={Jan de Leeuw},
            pdfborder={0 0 0},
            breaklinks=true}
\urlstyle{same}  % don't use monospace font for urls
\usepackage{color}
\usepackage{fancyvrb}
\newcommand{\VerbBar}{|}
\newcommand{\VERB}{\Verb[commandchars=\\\{\}]}
\DefineVerbatimEnvironment{Highlighting}{Verbatim}{commandchars=\\\{\}}
% Add ',fontsize=\small' for more characters per line
\usepackage{framed}
\definecolor{shadecolor}{RGB}{248,248,248}
\newenvironment{Shaded}{\begin{snugshade}}{\end{snugshade}}
\newcommand{\KeywordTok}[1]{\textcolor[rgb]{0.13,0.29,0.53}{\textbf{#1}}}
\newcommand{\DataTypeTok}[1]{\textcolor[rgb]{0.13,0.29,0.53}{#1}}
\newcommand{\DecValTok}[1]{\textcolor[rgb]{0.00,0.00,0.81}{#1}}
\newcommand{\BaseNTok}[1]{\textcolor[rgb]{0.00,0.00,0.81}{#1}}
\newcommand{\FloatTok}[1]{\textcolor[rgb]{0.00,0.00,0.81}{#1}}
\newcommand{\ConstantTok}[1]{\textcolor[rgb]{0.00,0.00,0.00}{#1}}
\newcommand{\CharTok}[1]{\textcolor[rgb]{0.31,0.60,0.02}{#1}}
\newcommand{\SpecialCharTok}[1]{\textcolor[rgb]{0.00,0.00,0.00}{#1}}
\newcommand{\StringTok}[1]{\textcolor[rgb]{0.31,0.60,0.02}{#1}}
\newcommand{\VerbatimStringTok}[1]{\textcolor[rgb]{0.31,0.60,0.02}{#1}}
\newcommand{\SpecialStringTok}[1]{\textcolor[rgb]{0.31,0.60,0.02}{#1}}
\newcommand{\ImportTok}[1]{#1}
\newcommand{\CommentTok}[1]{\textcolor[rgb]{0.56,0.35,0.01}{\textit{#1}}}
\newcommand{\DocumentationTok}[1]{\textcolor[rgb]{0.56,0.35,0.01}{\textbf{\textit{#1}}}}
\newcommand{\AnnotationTok}[1]{\textcolor[rgb]{0.56,0.35,0.01}{\textbf{\textit{#1}}}}
\newcommand{\CommentVarTok}[1]{\textcolor[rgb]{0.56,0.35,0.01}{\textbf{\textit{#1}}}}
\newcommand{\OtherTok}[1]{\textcolor[rgb]{0.56,0.35,0.01}{#1}}
\newcommand{\FunctionTok}[1]{\textcolor[rgb]{0.00,0.00,0.00}{#1}}
\newcommand{\VariableTok}[1]{\textcolor[rgb]{0.00,0.00,0.00}{#1}}
\newcommand{\ControlFlowTok}[1]{\textcolor[rgb]{0.13,0.29,0.53}{\textbf{#1}}}
\newcommand{\OperatorTok}[1]{\textcolor[rgb]{0.81,0.36,0.00}{\textbf{#1}}}
\newcommand{\BuiltInTok}[1]{#1}
\newcommand{\ExtensionTok}[1]{#1}
\newcommand{\PreprocessorTok}[1]{\textcolor[rgb]{0.56,0.35,0.01}{\textit{#1}}}
\newcommand{\AttributeTok}[1]{\textcolor[rgb]{0.77,0.63,0.00}{#1}}
\newcommand{\RegionMarkerTok}[1]{#1}
\newcommand{\InformationTok}[1]{\textcolor[rgb]{0.56,0.35,0.01}{\textbf{\textit{#1}}}}
\newcommand{\WarningTok}[1]{\textcolor[rgb]{0.56,0.35,0.01}{\textbf{\textit{#1}}}}
\newcommand{\AlertTok}[1]{\textcolor[rgb]{0.94,0.16,0.16}{#1}}
\newcommand{\ErrorTok}[1]{\textcolor[rgb]{0.64,0.00,0.00}{\textbf{#1}}}
\newcommand{\NormalTok}[1]{#1}
\usepackage{graphicx,grffile}
\makeatletter
\def\maxwidth{\ifdim\Gin@nat@width>\linewidth\linewidth\else\Gin@nat@width\fi}
\def\maxheight{\ifdim\Gin@nat@height>\textheight\textheight\else\Gin@nat@height\fi}
\makeatother
% Scale images if necessary, so that they will not overflow the page
% margins by default, and it is still possible to overwrite the defaults
% using explicit options in \includegraphics[width, height, ...]{}
\setkeys{Gin}{width=\maxwidth,height=\maxheight,keepaspectratio}
\IfFileExists{parskip.sty}{%
\usepackage{parskip}
}{% else
\setlength{\parindent}{0pt}
\setlength{\parskip}{6pt plus 2pt minus 1pt}
}
\setlength{\emergencystretch}{3em}  % prevent overfull lines
\providecommand{\tightlist}{%
  \setlength{\itemsep}{0pt}\setlength{\parskip}{0pt}}
\setcounter{secnumdepth}{5}
% Redefines (sub)paragraphs to behave more like sections
\ifx\paragraph\undefined\else
\let\oldparagraph\paragraph
\renewcommand{\paragraph}[1]{\oldparagraph{#1}\mbox{}}
\fi
\ifx\subparagraph\undefined\else
\let\oldsubparagraph\subparagraph
\renewcommand{\subparagraph}[1]{\oldsubparagraph{#1}\mbox{}}
\fi

%%% Use protect on footnotes to avoid problems with footnotes in titles
\let\rmarkdownfootnote\footnote%
\def\footnote{\protect\rmarkdownfootnote}

%%% Change title format to be more compact
\usepackage{titling}

% Create subtitle command for use in maketitle
\providecommand{\subtitle}[1]{
  \posttitle{
    \begin{center}\large#1\end{center}
    }
}

\setlength{\droptitle}{-2em}

  \title{Exceedingly Simple Principal Pivot Transforms}
    \pretitle{\vspace{\droptitle}\centering\huge}
  \posttitle{\par}
    \author{Jan de Leeuw}
    \preauthor{\centering\large\emph}
  \postauthor{\par}
      \predate{\centering\large\emph}
  \postdate{\par}
    \date{Version 001, March 23, 2016}


\begin{document}
\maketitle

{
\setcounter{tocdepth}{3}
\tableofcontents
}
Note: This is a working paper which will be expanded/updated frequently.
All suggestions for improvement are welcome. The directory
\href{http://deleeuwpdx.net/pubfolders/pivot}{deleeuwpdx.net/pubfolders/pivot}
has a pdf version, the complete Rmd file with all code chunks, the bib
file, and the R source code.

\section{Single Pivots}\label{single-pivots}

Suppose \(A\) is an \(n\times m\) matrix with \(a_{kk}\not= 0\). The
\(k^{th}\) \emph{principal single pivot} or \emph{PSP} transforms \(A\)
to the \(n\times m\) matrix \(\mathbf{piv}_k(A)\) with elements

\begin{equation}\label{E:pivot}
[\mathbf{piv}_k(A)]_{ij}=
\begin{cases}
\frac{1}{a_{kk}}&\text{ if }i=k\text{ and }j=k,\\
-\frac{a_{kj}}{a_{kk}}&\text{ if }i=k\text{ and }j\not=k,\\
\frac{a_{ik}}{a_{kk}}&\text{ if }i\not=k\text{ and }j=k,\\
a_{ij}-\frac{a_{ik}a_{kj}}{a_{kk}}&\text{ if }i\not=k\text{ and }j\not=k.
\end{cases}
\end{equation}

\(\mathbf{piv}_k(A)\) is undefined if \(a_{kk}=0\) . Note that the
matrix \(A\) is not necessarily square, symmetric, or positive
semi-definite.

If \(A\) is the matrix

\begin{verbatim}
##      [,1]  [,2]  [,3]  [,4]  [,5] 
## [1,]  1.00  1.00  1.00  1.00  1.00
## [2,]  1.00  2.00  2.00  2.00  2.00
## [3,]  1.00  2.00  3.00  3.00  3.00
## [4,]  1.00  2.00  3.00  4.00  4.00
## [5,]  1.00  2.00  3.00  4.00  5.00
\end{verbatim}

then pivoting on the second diagonal element, using our R function
\texttt{pivotOneRC()}, produces \(\mathbf{piv}_2(A)\)

\begin{verbatim}
## $a
##      [,1] [,2] [,3] [,4] [,5]
## [1,]  0.5  0.5    0    0    0
## [2,] -0.5  0.5   -1   -1   -1
## [3,]  0.0  1.0    1    1    1
## [4,]  0.0  1.0    1    2    2
## [5,]  0.0  1.0    1    2    3
## 
## $refuse
## [1] 0
\end{verbatim}

Note that the pivot operation is an involution, i.e.~it is its own
inverse, or \(\mathbf{piv}_k(\mathbf{piv}_k(A))=A\).

\begin{Shaded}
\begin{Highlighting}[]
\KeywordTok{mprint}\NormalTok{(}\KeywordTok{pivotOneRC}\NormalTok{(}\KeywordTok{pivotOneRC}\NormalTok{(a,}\DecValTok{2}\NormalTok{)}\OperatorTok{$}\NormalTok{a, }\DecValTok{2}\NormalTok{)}\OperatorTok{$}\NormalTok{a)}
\end{Highlighting}
\end{Shaded}

\begin{verbatim}
##      [,1]  [,2]  [,3]  [,4]  [,5] 
## [1,]  1.00  1.00  1.00  1.00  1.00
## [2,]  1.00  2.00  2.00  2.00  2.00
## [3,]  1.00  2.00  3.00  3.00  3.00
## [4,]  1.00  2.00  3.00  4.00  4.00
## [5,]  1.00  2.00  3.00  4.00  5.00
\end{verbatim}

PSP's are sometimes called \emph{sweeps}, and the \texttt{SWP()}
operator for symmetric matrices was introduced in statistical computing
by Beaton (1964), section 3.2. See Goodnight (1979), Stewart (1998) ,
p.~330--333, and Lange (2010), p.~95--99, for discussions of the sweep
operator in the context of numerical linear algebra. The version of the
sweep operator that became popular in statistics, however, is slightly
different from the one proposed by Beaton and from the one we use here.
It is

\begin{equation}\label{E:swp}
[\mathbf{swp}_k(A)]_{ij}=
\begin{cases}
-\frac{1}{a_{kk}}&\text{ if }i=k\text{ and }j=k,\\
\frac{a_{kj}}{a_{kk}}&\text{ if }i=k\text{ and }j\not=k,\\
\frac{a_{ik}}{a_{kk}}&\text{ if }i\not=k\text{ and }j=k,\\
a_{ij}-\frac{a_{ik}a_{kj}}{a_{kk}}&\text{ if }i\not=k\text{ and }j\not=k.
\end{cases}
\end{equation}

This version preserves symmetry, which is convenient because symmetry
can be used to minimize storage. But the symmetric sweep of
\$\eqref{E:swp} is not an involution. The inverse operation of this
sweep is the \emph{reverse sweep}, which is

\begin{equation}\label{E:rswp}
[\mathbf{rswp}_k(A)]_{ij}=
\begin{cases}
-\frac{1}{a_{kk}}&\text{ if }i=k\text{ and }j=k,\\
-\frac{a_{kj}}{a_{kk}}&\text{ if }i=k\text{ and }j\not=k,\\
-\frac{a_{ik}}{a_{kk}}&\text{ if }i\not=k\text{ and }j=k,\\
a_{ij}-\frac{a_{ik}a_{kj}}{a_{kk}}&\text{ if }i\not=k\text{ and }j\not=k.
\end{cases}
\end{equation}

In the funcion \texttt{pivotOneRC()} we have a type parameter. For type
equal to 1 we do a symmetric sweep, for type equal to 2 a reverse
symmetric sweep. Thus

\begin{Shaded}
\begin{Highlighting}[]
\KeywordTok{mprint}\NormalTok{ (}\KeywordTok{pivotOneRC}\NormalTok{(a, }\DecValTok{2}\NormalTok{, }\DataTypeTok{type =} \DecValTok{1}\NormalTok{)}\OperatorTok{$}\NormalTok{a)}
\end{Highlighting}
\end{Shaded}

\begin{verbatim}
##      [,1]  [,2]  [,3]  [,4]  [,5] 
## [1,]  0.50  0.50  0.00  0.00  0.00
## [2,]  0.50 -0.50  1.00  1.00  1.00
## [3,]  0.00  1.00  1.00  1.00  1.00
## [4,]  0.00  1.00  1.00  2.00  2.00
## [5,]  0.00  1.00  1.00  2.00  3.00
\end{verbatim}

\begin{Shaded}
\begin{Highlighting}[]
\KeywordTok{mprint}\NormalTok{ (}\KeywordTok{pivotOneRC}\NormalTok{(}\KeywordTok{pivotOneRC}\NormalTok{(a, }\DecValTok{2}\NormalTok{, }\DataTypeTok{type =} \DecValTok{1}\NormalTok{)}\OperatorTok{$}\NormalTok{a, }\DecValTok{2}\NormalTok{, }\DataTypeTok{type =} \DecValTok{2}\NormalTok{)}\OperatorTok{$}\NormalTok{a)}
\end{Highlighting}
\end{Shaded}

\begin{verbatim}
##      [,1]  [,2]  [,3]  [,4]  [,5] 
## [1,]  1.00  1.00  1.00  1.00  1.00
## [2,]  1.00  2.00  2.00  2.00  2.00
## [3,]  1.00  2.00  3.00  3.00  3.00
## [4,]  1.00  2.00  3.00  4.00  4.00
## [5,]  1.00  2.00  3.00  4.00  5.00
\end{verbatim}

Type equal to 3 is the PSP we started out with, and type equal to 4 is
its transpose, the involution

\begin{equation}\label{E:qiv}
[\mathbf{qiv}_k(A)]_{ij}=
\begin{cases}
\frac{1}{a_{kk}}&\text{ if }i=k\text{ and }j=k,\\
\frac{a_{kj}}{a_{kk}}&\text{ if }i=k\text{ and }j\not=k,\\
-\frac{a_{ik}}{a_{kk}}&\text{ if }i\not=k\text{ and }j=k,\\
a_{ij}-\frac{a_{ik}a_{kj}}{a_{kk}}&\text{ if }i\not=k\text{ and }j\not=k.
\end{cases}
\end{equation}

The default in \texttt{pivotOneRC()} is to set type equal to 3, and we
will use this default throughout the rest of the paper.

\section{Sequences of pivots}\label{sequences-of-pivots}

Our R function \texttt{pivotRC()} performs a sequence of pivots on a
matrix. If both PSP's are defined then
\(\mathbf{piv}_k(\mathbf{piv}_\ell(A))=\mathbf{piv}_\ell(\mathbf{piv}_k(A))\),
and thus the results are independent of the order in which we pivot. We
obtain numerical stability by always pivoting on the largest diagonal
element in the remainder of the sequence of pivots. This also ensures
the function gives a sensible result if the matrices involved are
singular.

Consider the matrix we get from our previous one by setting element
\((1,1)\) to zero. The resulting matrix is still non-singular.

\begin{verbatim}
##      [,1]  [,2]  [,3]  [,4]  [,5] 
## [1,]  0.00  1.00  1.00  1.00  1.00
## [2,]  1.00  2.00  2.00  2.00  2.00
## [3,]  1.00  2.00  3.00  3.00  3.00
## [4,]  1.00  2.00  3.00  4.00  4.00
## [5,]  1.00  2.00  3.00  4.00  5.00
\end{verbatim}

If we pivot on the first four elements we get

\begin{verbatim}
##      [,1]  [,2]  [,3]  [,4]  [,5] 
## [1,] -2.00  1.00  0.00  0.00  0.00
## [2,]  1.00  1.00 -1.00  0.00  0.00
## [3,]  0.00 -1.00  2.00 -1.00 -0.00
## [4,] -0.00  0.00 -1.00  1.00 -1.00
## [5,]  0.00  0.00  0.00  1.00  1.00
\end{verbatim}

and pivots are done in the order 4, 2, 1, 3.

Another example is the rank 2 matrix

\begin{Shaded}
\begin{Highlighting}[]
\KeywordTok{mprint}\NormalTok{(a<-}\KeywordTok{tcrossprod}\NormalTok{(}\KeywordTok{matrix}\NormalTok{(}\KeywordTok{c}\NormalTok{(}\DecValTok{1}\NormalTok{,}\DecValTok{1}\NormalTok{,}\DecValTok{1}\NormalTok{,}\DecValTok{1}\NormalTok{,}\DecValTok{1}\NormalTok{,}\OperatorTok{-}\DecValTok{1}\NormalTok{,}\OperatorTok{-}\DecValTok{1}\NormalTok{,}\DecValTok{1}\NormalTok{),}\DecValTok{4}\NormalTok{,}\DecValTok{2}\NormalTok{))}\OperatorTok{/}\DecValTok{2}\NormalTok{)}
\end{Highlighting}
\end{Shaded}

\begin{verbatim}
##      [,1]  [,2]  [,3]  [,4] 
## [1,]  1.00  0.00  0.00  1.00
## [2,]  0.00  1.00  1.00  0.00
## [3,]  0.00  1.00  1.00  0.00
## [4,]  1.00  0.00  0.00  1.00
\end{verbatim}

A complete pivot uses the order 1, 2, 3, 4 and finds

\begin{Shaded}
\begin{Highlighting}[]
\KeywordTok{mprint}\NormalTok{(h}\OperatorTok{$}\NormalTok{a)}
\end{Highlighting}
\end{Shaded}

\begin{verbatim}
##      [,1]  [,2]  [,3]  [,4] 
## [1,]  1.00 -0.00  0.00 -1.00
## [2,] -0.00  1.00 -1.00 -0.00
## [3,]  0.00  1.00  0.00  0.00
## [4,]  1.00  0.00  0.00  0.00
\end{verbatim}

Note that if the function \texttt{pivotOneRC()} is asked to pivot on a
zero diagonal element, it refuses and just returns its argument. Thus if
zero diagonal elements are encountered, it is just as if the
corresponding indices are not in the sequence, and we actually pivot on
a shorter sequence. \texttt{pivotRC()} returns an indicator of the
pivots that are skipped, which for our last example is 0, 0, 1, 1.

Also note that order-invariance for \texttt{pivotRC()} is no longer true
if pivots are skipped. Consider

\begin{verbatim}
##      [,1]  [,2]  [,3]  [,4] 
## [1,]  0.00  0.00  0.00  1.00
## [2,]  0.00  0.00  0.00  1.00
## [3,]  0.00  0.00  0.00  1.00
## [4,]  1.00  1.00  1.00  1.00
\end{verbatim}

Pivoting over all indices, but in a different order, gives different
results.

\begin{Shaded}
\begin{Highlighting}[]
\KeywordTok{pivotRC}\NormalTok{(a, }\DecValTok{1}\OperatorTok{:}\DecValTok{4}\NormalTok{)}
\end{Highlighting}
\end{Shaded}

\begin{verbatim}
## $a
##      [,1] [,2] [,3] [,4]
## [1,]   -1   -1   -1    1
## [2,]    1    0    0    0
## [3,]    1    0    0    0
## [4,]    1    0    0    0
## 
## $pivot
## [1] 4 1 2 3
## 
## $skip
## [1] 0 0 1 1
\end{verbatim}

\begin{Shaded}
\begin{Highlighting}[]
\KeywordTok{pivotRC}\NormalTok{(a, }\DecValTok{4}\OperatorTok{:}\DecValTok{1}\NormalTok{)}
\end{Highlighting}
\end{Shaded}

\begin{verbatim}
## $a
##      [,1] [,2] [,3] [,4]
## [1,]    0    0    1    0
## [2,]    0    0    1    0
## [3,]   -1   -1   -1    1
## [4,]    0    0    1    0
## 
## $pivot
## [1] 4 3 2 1
## 
## $skip
## [1] 0 0 1 1
\end{verbatim}

In the first case we compute \(\mathbf{piv}_{\{1,4\}}(A)\), in the
second case \(\mathbf{piv}_{\{3,4\}}(A)\), and these are different.

Suppose

\begin{equation}
A=\begin{bmatrix}\alpha&a'\\b&C\end{bmatrix},\label{E:psd}
\end{equation}

with \(\alpha\not= 0\). Then \[
\mathbf{piv}_1(A)=\begin{bmatrix}\frac{1}{\alpha}&\frac{-a'}{\alpha}\\\frac{b}{\alpha}&C-\frac{ba'}{\alpha}\end{bmatrix}.
\] The Schur complement (Ouellette (1981), Zhang (2005)) of
\(\frac{1}{\alpha}\) in \(\mathbf{piv}_1(A)\) is \(C\). Since rank is
additive over the Schur complement, we have \[
\mathbf{rank}(\mathbf{piv}_1(A))=1+\mathbf{rank}(C),
\] and thus \(\mathbf{piv}_1(A)\) is non-singular whenever \(A\) is.

For a symmetric \(A\) in \(\eqref{E:psd}\) we have that \(A\) is
positive definite if and only if \(\alpha>0\) and
\(C-\frac{ba'}{\alpha}\) is positive definite (Bekker (1988)). Thus if
\(A\) is positive definite, all subsequent pivot elements are positive.
If \(A\) is unit triangular, either upper or lower, then
\(C-\frac{ba'}{\alpha}\) is unit triangular as well, and thus all
subsequent pivot elements are equal to one.

Of course the choice of pivot \(k=1\) in \(\eqref{E:psd}\) was for
notational convenience only, the results are true for any \(k\).

\section{The partial inverse}\label{the-partial-inverse}

Suppose \(K\) is a subset of \(N:=\{1,2,\cdots,n\}\) and
\(\underline{K}\) is the complimentary subset. In our examples we
usually take, for convenience, \(K=\{1,2,\cdots,k\}\). The
\emph{principal pivot transform} or \emph{PPT} of \(A\) on \(K\), which
generalizes the simple pivot, is \[
\mathbf{piv}_K(A)=\begin{bmatrix}
A_{KK}^{-1}&-A_{KK}^{-1}A_{K\underline{K}}\\
A_{\underline K K}A_{KK^{}}^{-1}&A_{\underline K\underline K}-A_{\underline K K}A_{KK^{}}^{-1}A_{K\underline{K}}
\end{bmatrix}.
\] The PPT is defined only if \(A_{KK}\) is non-singular. A
\emph{complete} PPT is defined as \(\mathbf{piv}_K()\) for \(K=N\).

The PPT of a square matrix \(A\) is called the \emph{partial inverse} by
Wermuth, Wiedenbeck, and Cox (2006) and Wiedenbeck and Wermuth (2010).
The name makes sense because if \[
\begin{bmatrix}
A_{KK}&A_{\underline K}\\
A_{\underline K K}&A_{\underline K\underline K}
\end{bmatrix}
\begin{bmatrix}
x\\y
\end{bmatrix}=
\begin{bmatrix}
u\\v
\end{bmatrix},
\] then \[
\begin{bmatrix}
A_{KK}^{-1}&-A_{KK}^{-1}A_{K\underline{K}}\\
A_{\underline K K}A_{KK^{}}^{-1}&A_{\underline K\underline K}-A_{\underline K K}A_{KK^{}}^{-1}A_{K\underline{K}}
\end{bmatrix}
\begin{bmatrix}
u\\y
\end{bmatrix}=
\begin{bmatrix}
x\\v
\end{bmatrix},
\] And conversely. Of course all this assumes that \(A_{KK}\) is
non-singular.

The following results are true for arbitrary index sets \(K\) and \(L\)
and their complement \(\underline K\) and \(\underline L\). For a proof,
see Wermuth, Wiedenbeck, and Cox (2006), lemma 2.1. They assume that all
principal submatrices are invertible, which is true for positive
definite matrices and unit triangular matrices. In that case all PPT's
are defined.

\begin{itemize}
\tightlist
\item
  \(A=\mathbf{piv}_K(\mathbf{piv}_K(A))\)
\item
  \(\mathbf{piv}_K(A)=(\mathbf{piv}_{\underline K}(A))^{-1}\)
\item
  \(\mathbf{piv}_{\underline K}(\mathbf{piv}_K(A))=A^{-1}\)
\item
  \(\mathbf{piv}_L(\mathbf{piv}_K(A))=\mathbf{piv}_K(\mathbf{piv}_L(A))\)
\item
  \(\mathbf{piv}_K(A)=\mathbf{piv}_{\underline K}(A^{-1})\)
\item
  \((\mathbf{piv}_K(A))^{-1}=\mathbf{piv}_K(A^{-1})\)
\end{itemize}

Wermuth, Wiedenbeck, and Cox (2006) summarize the main properties of
partial inversion in their theorem 2.2. Suppose \(K,L,\) and \(M\)
partition the \(n\) indices of a square matrix, and the PPT's are
defined Then

\begin{itemize}
\tightlist
\item
  Commutativity:
  \(\mathbf{piv}_K(\mathbf{piv}_L(A))=\mathbf{piv}_K(\mathbf{piv}_L(A))=\mathbf{piv}_H(A)\).
\item
  Exchangeability:
  \([\mathbf{piv}_K(A)]_{K\cup L,K\cup L}=\mathbf{piv}_K(A_{K\cup L,K\cup L})\).
\item
  Symmetric Difference:
  \(\mathbf{piv}_{K\cup L}(A)\mathbf{piv}_{L\cup M}(A)=\mathbf{piv}_{K\cup M}(A)\).
\end{itemize}

If \(A\) is not square, then some of these results still make sense. For
both a \emph{broad} \(n\times m\) matrix A, with \(n<m\), and a
\emph{tall} \(A\), with \(n>m\), we can pivot on the diagonal elements
and the defining equations for the principal pivot transform still
apply.

Note that a PPT can be computed by a sequence of PSP's, provided all
PSP's on the way are defined. But consider the example

\begin{verbatim}
##      [,1]  [,2]  [,3]  [,4] 
## [1,]  0.00  1.00  1.00  0.00
## [2,]  1.00  0.00  0.00  1.00
## [3,]  1.00  0.00  1.00  0.00
## [4,]  0.00  1.00  0.00  1.00
\end{verbatim}

Now \(\mathbf{piv}_1(A)\) and \(\mathbf{piv}_2(A)\) are both not
defined, and thus our function \texttt{pivotRC(a,\ 1:2)} does nothing,
while \[
\mathbf{piv}_{\{1,2\}}(A)=\begin{bmatrix}
\hphantom{+}0&\hphantom{+}1&\hphantom{+}0&-1\\
\hphantom{+}1&\hphantom{+}0&-1&\hphantom{+}0\\
\hphantom{+}0&\hphantom{+}1&\hphantom{+}1&-1\\
\hphantom{+}1&\hphantom{+}0&-1&\hphantom{+}1
\end{bmatrix}.
\]

\section{Inverse and determinant}\label{inverse-and-determinant}

For a square matrix \[
A=\begin{bmatrix}\alpha&a'\\b&C\end{bmatrix},
\] with \(\alpha\not= 0\), by Schur's determinant formula, \[
\mathbf{det}(A)=\alpha\ \mathbf{det}(C-\frac{ba'}{\alpha}),
\] and thus \(\mathbf{det}(A)\) is the product of the pivots in a
complete PPT.

If \(A\) is square and non-singular, then carrying out a complete PPT
produces the inverse of the matrix, or \(\mathbf{piv}_N(A)=A^{-1}\).
This assumes, of course, that a complete PPT can actually be carried out
and we do not run into a zero pivot.

If \(A\) is singular, say of rank \(r\), then we can permute rows and
columns so that we have the leading principal submatrix of order \(r\)
is non-singular. This easily follows from the pivoted LU decomposition,
see Stewart (1998), theorem 2.13. If \[
A=\begin{bmatrix}B&C\\D&E\end{bmatrix}
\] with \(B\) an \(r\times r\) matrix such that
\(\mathbf{rank}(B)=\mathbf{rank}(A)=r\), and a pivot of \(A\) over the
first \(r\) elements is possible, then \[
\mathbf{piv}_{\{1,2,\cdots,r\}}(A)=\begin{bmatrix}B^{-1}&-B^{-1}C\\DB^{-1}&0\end{bmatrix}.
\] This implies that
\(A\left(\mathbf{piv}_{\{1,2,\cdots,r\}}(A)\right)A=A\), i.e.
\(\mathbf{piv}_{\{1,2,\cdots,r\}}(A)\) satisfies the first Penrose
condition, i.e.~is a (1)-inverse of \(A\).

\section{Least Squares}\label{least-squares}

Statisticians like the \(\mathbf{piv}()\) and \(\mathbf{swp}()\)
operators, because of their connection with linear least squares
regression, and with stepwise regression in particular (Jennrich
(1977)). In addition there are connections with covariance selection
(Dempster (1972)) and with graphical models (Wermuth, Wiedenbeck, and
Cox (2006)).

For the least squares problem of minimizing \((y-Xb)'(y-Xb)\) the normal
equations are \(X'Xb=X'y\). Suppose \(X=(X_1\mid X_2)\), with \(m_1\)
and \(m_2\) columns, and consider the matrix \[
A=\begin{bmatrix}
X_1'X_1&X_1'X_2&X_1'y\\
X_2'X_1&X_2'X_2&X_2'y\\
y'X_1&y'X_2&y'y
\end{bmatrix}
\] Pivoting on the \(m_1\) elements in \(X_1'X_1\), assuming it is
invertible, gives \[
\mathbf{piv}_{\{1,2,\cdots,m\}}(A)=
\begin{bmatrix}
(X_1'X_1)^{-1}&-(X_1'X_1)^{-1}X_1'X_2&-(X_1'X_1)^{-1}X_1'y\\
X_2'X_1(X_1'X_1)^{-1}&X_2'(I-X_1(X_1'X_1)^{-1}X_1')X_2&X_2'(I-X_1(X_1'X_1)^{-1}X_1')y\\
y'X_1(X_1'X_1)^{-1}&y'(I-X_1(X_1'X_1)^{-1}X_1')X_2&y'(I-X_1(X_1'X_1)^{-1}X_1')y
\end{bmatrix}.
\] All submatrices in the PPT give statistics which are relevant in the
context of the least squares problem. We have the regression
coefficients, their dispersion matrix, the residual sum of squares, and
the partial covariances of \(X_2\) given \(X_1\). Moreover an additional
PSP on one of the \(m_2\) elements from \(X_2\) will add a variable to
the regression, and repeating a PSP from the first \(m_1\) will remove
that variable from the regression. \#Code

The basic PSP routine \texttt{pivotOneC()} is written in C. If a pivot
is zero, it is skipped. Performing a number of pivots to get a PPT is
done in \texttt{pivotC()}, also in C. The \texttt{.C()} interface is
used in R to write wrappers \texttt{pivotOneRC()} and \texttt{pivotRC()}
for the C routines. Note that the \texttt{ggm} package from CRAN
(Marchetti, Drton, and Sadeghi (2015)) has a \texttt{swp()} function
that uses \texttt{solve()} from base R to compute the partial inverse.
Pivoting on principal submatrices to compute the partial inverse will
generally be faster than building up the partial inverse from pivoting
on a sequence of diagonal elements like we do.

In previous work (({\textbf{???}})) we implemented sweeping by writing a
simple R wrapper for the fortran routine in AS 178 (Clarke (1982)). We
improve this by switching to principal pivoting, to C, and to
incorporating reordering the remaining pivots for numerical stability,
similar to the suggestions of Ridout and Cobby (1989).

\subsection{R Code}\label{r-code}

\begin{Shaded}
\begin{Highlighting}[]
\KeywordTok{dyn.load}\NormalTok{ (}\StringTok{"pivot.so"}\NormalTok{)}

\NormalTok{pivotOneRC <-}\StringTok{ }\ControlFlowTok{function}\NormalTok{ (a, k, }\DataTypeTok{type =} \DecValTok{3}\NormalTok{, }\DataTypeTok{eps =} \FloatTok{1e-10}\NormalTok{) \{}
\NormalTok{  n <-}\StringTok{ }\KeywordTok{nrow}\NormalTok{ (a)}
\NormalTok{  m <-}\StringTok{ }\KeywordTok{ncol}\NormalTok{ (a)}
\NormalTok{  h <-}
\StringTok{    }\KeywordTok{.C}\NormalTok{(}
      \StringTok{"pivotOneC"}\NormalTok{,}
      \DataTypeTok{a =} \KeywordTok{as.double}\NormalTok{ (a),}
      \DataTypeTok{k =} \KeywordTok{as.integer}\NormalTok{(k),}
      \DataTypeTok{n =} \KeywordTok{as.integer}\NormalTok{ (n),}
      \DataTypeTok{m =} \KeywordTok{as.integer}\NormalTok{ (m),}
      \DataTypeTok{type =} \KeywordTok{as.integer}\NormalTok{ (type),}
      \DataTypeTok{refuse =} \KeywordTok{as.integer}\NormalTok{ (}\DecValTok{0}\NormalTok{),}
      \DataTypeTok{eos =} \KeywordTok{as.double}\NormalTok{ (eps)}
\NormalTok{    )}
  \KeywordTok{return}\NormalTok{ (}\KeywordTok{list}\NormalTok{(}\DataTypeTok{a =} \KeywordTok{matrix}\NormalTok{(h}\OperatorTok{$}\NormalTok{a, n, m), }\DataTypeTok{refuse =}\NormalTok{ h}\OperatorTok{$}\NormalTok{refuse))}
\NormalTok{\}}

\NormalTok{pivotRC <-}\StringTok{ }\ControlFlowTok{function}\NormalTok{ (a, ind, }\DataTypeTok{type =} \DecValTok{3}\NormalTok{, }\DataTypeTok{eps =} \FloatTok{1e-10}\NormalTok{) \{}
\NormalTok{  n <-}\StringTok{ }\KeywordTok{nrow}\NormalTok{ (a)}
\NormalTok{  m <-}\StringTok{ }\KeywordTok{ncol}\NormalTok{ (a)}
\NormalTok{  p <-}\StringTok{ }\KeywordTok{length}\NormalTok{ (ind)}
\NormalTok{  done <-}\StringTok{ }\KeywordTok{rep}\NormalTok{ (}\DecValTok{0}\NormalTok{, p)}
\NormalTok{  h <-}
\StringTok{    }\KeywordTok{.C}\NormalTok{(}
      \StringTok{"pivotC"}\NormalTok{,}
      \DataTypeTok{a =} \KeywordTok{as.double}\NormalTok{ (a),}
      \DataTypeTok{ind =} \KeywordTok{as.integer}\NormalTok{ (ind),}
      \DataTypeTok{jnd =} \KeywordTok{as.integer}\NormalTok{ (}\KeywordTok{rep}\NormalTok{ (}\DecValTok{0}\NormalTok{, p)),}
      \DataTypeTok{done =} \KeywordTok{as.integer}\NormalTok{ (}\KeywordTok{rep}\NormalTok{ (}\DecValTok{0}\NormalTok{, p)),}
      \DataTypeTok{n =} \KeywordTok{as.integer}\NormalTok{ (n),}
      \DataTypeTok{m =} \KeywordTok{as.integer}\NormalTok{ (m),}
      \DataTypeTok{p =} \KeywordTok{as.integer}\NormalTok{ (p),}
      \DataTypeTok{type =} \KeywordTok{as.integer}\NormalTok{ (type),}
      \DataTypeTok{skip =} \KeywordTok{as.integer}\NormalTok{ (}\KeywordTok{rep}\NormalTok{ (}\DecValTok{0}\NormalTok{, p)),}
      \DataTypeTok{eps =} \KeywordTok{as.double}\NormalTok{ (eps)}
\NormalTok{    )}
  \KeywordTok{return}\NormalTok{ (}\KeywordTok{list}\NormalTok{(}\DataTypeTok{a =} \KeywordTok{matrix}\NormalTok{(h}\OperatorTok{$}\NormalTok{a, n, m), }\DataTypeTok{pivot =}\NormalTok{ h}\OperatorTok{$}\NormalTok{jnd, }\DataTypeTok{skip =}\NormalTok{ h}\OperatorTok{$}\NormalTok{skip))}
\NormalTok{\}}
\end{Highlighting}
\end{Shaded}

\subsection{C Code}\label{c-code}

\begin{Shaded}
\begin{Highlighting}[]
\DataTypeTok{void}\NormalTok{ pivotOneC (}\DataTypeTok{double}\NormalTok{ *, }\DataTypeTok{int}\NormalTok{ *, }\DataTypeTok{int}\NormalTok{ *, }\DataTypeTok{int}\NormalTok{ *, }\DataTypeTok{int}\NormalTok{ *, }\DataTypeTok{int}\NormalTok{ *, }\DataTypeTok{double}\NormalTok{ *eps);}
\DataTypeTok{void}\NormalTok{ pivotC (}\DataTypeTok{double}\NormalTok{ *, }\DataTypeTok{int}\NormalTok{ *, }\DataTypeTok{int}\NormalTok{ *, }\DataTypeTok{int}\NormalTok{ *, }\DataTypeTok{int}\NormalTok{ *, }\DataTypeTok{int}\NormalTok{ *, }\DataTypeTok{int}\NormalTok{ *, }\DataTypeTok{int}\NormalTok{ *, }\DataTypeTok{int}\NormalTok{ *, }\DataTypeTok{double}\NormalTok{ *eps);}
\DataTypeTok{void}\NormalTok{ mprint (}\DataTypeTok{double}\NormalTok{ *, }\DataTypeTok{int}\NormalTok{ *, }\DataTypeTok{int}\NormalTok{ *);}

\DataTypeTok{void}\NormalTok{ pivotOneC (}\DataTypeTok{double}\NormalTok{ *a, }\DataTypeTok{int}\NormalTok{ *kk, }\DataTypeTok{int}\NormalTok{ *nn, }\DataTypeTok{int}\NormalTok{ *mm, }\DataTypeTok{int}\NormalTok{ *type, }\DataTypeTok{int}\NormalTok{ *refuse, }\DataTypeTok{double}\NormalTok{ *eps) \{}
    \DataTypeTok{int}\NormalTok{ i, j, ik, kj, ij, kp, s00, s01, s10, n = *nn, m = *mm, k = *kk, tp = *type;}
    \DataTypeTok{double}\NormalTok{ pv, ee = *eps;}
\NormalTok{    *refuse = }\DecValTok{0}\NormalTok{;}
    \ControlFlowTok{if}\NormalTok{ (tp == }\DecValTok{1}\NormalTok{) \{}
\NormalTok{        s00 = }\DecValTok{-1}\NormalTok{;}
\NormalTok{        s01 = }\DecValTok{1}\NormalTok{;}
\NormalTok{        s10 = }\DecValTok{1}\NormalTok{;}
\NormalTok{    \}}
    \ControlFlowTok{if}\NormalTok{ (tp == }\DecValTok{2}\NormalTok{) \{}
\NormalTok{        s00 = }\DecValTok{-1}\NormalTok{;}
\NormalTok{        s01 = }\DecValTok{-1}\NormalTok{;}
\NormalTok{        s10 = }\DecValTok{-1}\NormalTok{;}
\NormalTok{    \}}
   \ControlFlowTok{if}\NormalTok{ (tp == }\DecValTok{3}\NormalTok{) \{}
\NormalTok{        s00 = }\DecValTok{1}\NormalTok{;}
\NormalTok{        s01 = }\DecValTok{-1}\NormalTok{;}
\NormalTok{        s10 = }\DecValTok{1}\NormalTok{;}
\NormalTok{    \}}
   \ControlFlowTok{if}\NormalTok{ (tp == }\DecValTok{4}\NormalTok{) \{}
\NormalTok{        s00 = }\DecValTok{1}\NormalTok{;}
\NormalTok{        s01 = }\DecValTok{1}\NormalTok{;}
\NormalTok{        s10 = }\DecValTok{-1}\NormalTok{;}
\NormalTok{    \}}
\NormalTok{    kp = MINDEX(k, k, n);}
\NormalTok{    pv = a[kp];}
    \ControlFlowTok{if}\NormalTok{ (fabs (pv) < ee) \{}
\NormalTok{        *refuse = }\DecValTok{1}\NormalTok{; }
        \ControlFlowTok{return}\NormalTok{;}
\NormalTok{    \}}
    \ControlFlowTok{for}\NormalTok{ (i = }\DecValTok{1}\NormalTok{; i <= n; i++) \{}
        \ControlFlowTok{if}\NormalTok{ (i == k) }\ControlFlowTok{continue}\NormalTok{;}
\NormalTok{        ik = MINDEX (i, k , n);}
        \ControlFlowTok{for}\NormalTok{ (j = }\DecValTok{1}\NormalTok{; j <= m; j++) \{}
            \ControlFlowTok{if}\NormalTok{ (j == k) }\ControlFlowTok{continue}\NormalTok{;}
\NormalTok{            kj = MINDEX (k, j, n);}
\NormalTok{            ij = MINDEX (i, j, n);}
\NormalTok{            a[ij] = a[ij] - a[ik] * a[kj] / pv;}
\NormalTok{        \}}
\NormalTok{    \}}
    \ControlFlowTok{for}\NormalTok{ (i = }\DecValTok{1}\NormalTok{; i <= n; i++) \{}
        \ControlFlowTok{if}\NormalTok{ (i == k) }\ControlFlowTok{continue}\NormalTok{;}
\NormalTok{        ik = MINDEX (i, k, n);}
\NormalTok{        a[ik] = s10 * a[ik] / pv;}
\NormalTok{    \}}
    \ControlFlowTok{for}\NormalTok{ (j = }\DecValTok{1}\NormalTok{; j <= m; j++) \{}
        \ControlFlowTok{if}\NormalTok{ (j == k) }\ControlFlowTok{continue}\NormalTok{;}
\NormalTok{        kj = MINDEX (k, j, n);}
\NormalTok{        a[kj] = s01 * a[kj] / pv;}
\NormalTok{    \}}
\NormalTok{    a[kp] = s00 / pv;}
\NormalTok{\}}


\DataTypeTok{void}\NormalTok{ pivotC (}\DataTypeTok{double}\NormalTok{ *a, }\DataTypeTok{int}\NormalTok{ *ind, }\DataTypeTok{int}\NormalTok{ *jnd, }\DataTypeTok{int}\NormalTok{ *done, }\DataTypeTok{int}\NormalTok{ *nn, }\DataTypeTok{int}\NormalTok{ *mm, }\DataTypeTok{int}\NormalTok{ *pp, }\DataTypeTok{int}\NormalTok{ *type, }\DataTypeTok{int}\NormalTok{ *skip, }\DataTypeTok{double}\NormalTok{ *eps) \{}
    \DataTypeTok{int}\NormalTok{ ipiv, kpiv, lpiv, refuse;}
    \DataTypeTok{int}\NormalTok{ n = *nn, m = *mm, p = *pp;}
    \DataTypeTok{double}\NormalTok{ fmax, fpiv;}
    \ControlFlowTok{for}\NormalTok{ (}\DataTypeTok{int}\NormalTok{ l = }\DecValTok{0}\NormalTok{; l < p; l++) \{}
\NormalTok{        jnd[l] = }\DecValTok{0}\NormalTok{;}
\NormalTok{        done[l] = }\DecValTok{0}\NormalTok{;}
\NormalTok{        skip[l] = }\DecValTok{0}\NormalTok{;}
\NormalTok{    \}}
    \ControlFlowTok{for}\NormalTok{ (}\DataTypeTok{int}\NormalTok{ l = }\DecValTok{1}\NormalTok{; l <= p; l++) \{}
\NormalTok{        fmax = }\DecValTok{-1}\NormalTok{.}\DecValTok{0}\NormalTok{;}
        \ControlFlowTok{for}\NormalTok{ (}\DataTypeTok{int}\NormalTok{ k = }\DecValTok{1}\NormalTok{; k <= p; k++) \{}
\NormalTok{            kpiv = ind[k - }\DecValTok{1}\NormalTok{];}
            \ControlFlowTok{if}\NormalTok{ (done[k - }\DecValTok{1}\NormalTok{] == }\DecValTok{1}\NormalTok{) }\ControlFlowTok{continue}\NormalTok{;}
\NormalTok{            fpiv = fabs (a[MINDEX (kpiv, kpiv, n)]);}
            \ControlFlowTok{if}\NormalTok{ (fpiv > fmax) \{}
\NormalTok{                ipiv = k;}
\NormalTok{                lpiv = kpiv;}
\NormalTok{                fmax = fpiv;}
\NormalTok{            \}}
\NormalTok{        \}}
\NormalTok{        done[ipiv - }\DecValTok{1}\NormalTok{] = }\DecValTok{1}\NormalTok{;}
\NormalTok{        jnd [l - }\DecValTok{1}\NormalTok{] = lpiv;}
\NormalTok{        pivotOneC(a, &lpiv, &n, &m, type, &refuse, eps);}
\NormalTok{        skip[l - }\DecValTok{1}\NormalTok{] = refuse;}
\NormalTok{    \}}
\NormalTok{\}}


\DataTypeTok{void}\NormalTok{ mprint (}\DataTypeTok{double}\NormalTok{ * a, }\DataTypeTok{int}\NormalTok{ *nn, }\DataTypeTok{int}\NormalTok{ *mm) \{}
    \DataTypeTok{int}\NormalTok{ i, j, ij, n = *nn, m = *mm;}
    \ControlFlowTok{for}\NormalTok{ (i = }\DecValTok{1}\NormalTok{; i <= n; i++) \{}
        \ControlFlowTok{for}\NormalTok{ (j = }\DecValTok{1}\NormalTok{; j <= m; j++) \{}
\NormalTok{            ij = MINDEX (i, j, n);}
\NormalTok{            printf (}\StringTok{" }\SpecialCharTok{%5.3f}\StringTok{ "}\NormalTok{, a[ij]);}
\NormalTok{        \}}
\NormalTok{    printf (}\StringTok{"}\SpecialCharTok{\textbackslash{}n}\StringTok{"}\NormalTok{);}
\NormalTok{    \}}
\NormalTok{\}}
\end{Highlighting}
\end{Shaded}

\section{NEWS}\label{news}

\section*{References}\label{references}
\addcontentsline{toc}{section}{References}

\hypertarget{refs}{}
\hypertarget{ref-beaton_64}{}
Beaton, A.E. 1964. ``The Use of Special Matrix Operators in Statistical
Calculus.'' RB 64-51. Princeton, NJ: Educational Testing Service.

\hypertarget{ref-bekker_88}{}
Bekker, P.A. 1988. ``The Positive Semidefiniteness of Partitioned
Matrices.'' \emph{Linear Algebra and Its Applications} 111: 261--78.

\hypertarget{ref-clarke_82}{}
Clarke, M.R.B. 1982. ``Algorithm AS 178: The Gauss-Jordan Sweep Operator
with Detection of Collinearity.'' \emph{Journal of the Royal Statistical
Society. Series C (Applied Statistics),} 31: 166--68.

\hypertarget{ref-dempster_72}{}
Dempster, A.P. 1972. ``Covariance Selection.'' \emph{Biometrics} 28:
157--75.

\hypertarget{ref-goodnight_79}{}
Goodnight, J.H. 1979. ``A Tutorial on the SWEEP Operator.''
\emph{American Statistician} 33: 149--58.

\hypertarget{ref-jennrich_77}{}
Jennrich, R.I. 1977. ``Stepwise Regression.'' In \emph{Statistical
Methods for Digital Computers}, edited by Enslein K., A. Ralston, and
H.S. Wilf. Wiley.

\hypertarget{ref-lange_10}{}
Lange, K. 2010. \emph{Numerical Analysis for Statisticians}. Second
Edition. Springer.

\hypertarget{ref-marchetti_drton_sadeghi_15}{}
Marchetti, G.M., M. Drton, and K. Sadeghi. 2015. \emph{ggm: Functions
for Graphical Markov Models}.
\url{https://CRAN.R-project.org/package=ggm}.

\hypertarget{ref-ouellette_81}{}
Ouellette, D.V. 1981. ``Schur Complements and Statistics.'' \emph{Linear
Algebra and Its Applications} 36: 187--295.

\hypertarget{ref-ridout_cobby_89}{}
Ridout, M.S., and J. M. Cobby. 1989. ``Remark AS R78: A Remark on
Algorithm AS 178: The Gauss-Jordan Sweep Operator with Detection of
Collinearity.'' \emph{Journal of the Royal Statistical Society. Series C
(Applied Statistics),} 38: 420--22.

\hypertarget{ref-stewart_98}{}
Stewart, G.W. 1998. \emph{Matrix Algorithms. Volume I: Basic
Decompositions}. SIAM.

\hypertarget{ref-wermuth_wiedenbeck_cox_06}{}
Wermuth, N., M. Wiedenbeck, and D.R. Cox. 2006. ``Partial Inversion for
Linear Systems and Partial Closures of Independence Graphs.'' \emph{BIT
Numerical Mathematics} 46: 883--901.

\hypertarget{ref-wiedenbeck_wermuth_10}{}
Wiedenbeck, M., and N. Wermuth. 2010. ``Changing Parameters by Partial
Mappings.'' \emph{Statistica Sinica} 20: 823--36.

\hypertarget{ref-zhang_05}{}
Zhang, F., ed. 2005. \emph{The Schur Complement and its Applications}.
Vol. 4. Numerical Methods and Algorithms. Springer.


\end{document}
